\section{\label{sec:intro}Introduction}
Organizations wishing to train LLMs on sensitive records are expected to do so under a formal privacy guarantee.
The standard answers are private optimization by differentially private stochastic gradient descent (DP-SGD),
which adds calibrated noise to the gradient updates \citep{abadi2016deep, mironov2017renyi},
and text sanitization,
altering the training data before the model sees it \citep{feyisetan2020privacy, yue2021differential, carvalho2021tem}.
Both are mature enough to deploy,
and both are routinely evaluated the same way:
fit at a range of budgets,
report benchmark accuracy,
and read the cost of privacy off the gap to a non-private baseline.

\textbf{Accuracy is the wrong instrument.}
What privacy costs shows up in how the model behaves,
beyond how often it is right.
An adapted model can keep its accuracy and still follow instructions less reliably,
fabricate the facts it used to ground,
or stop refusing requests it should decline;
we call this alignment degradation.
That private training induces it is established \citep{ramesh2026privacyhallucination},
with the damage falling on facts learned during fine-tuning rather than on what the model already knew.
The cause is open,
and it decides the fix.
If noise destabilizes optimization,
the fix is a better optimizer.
If the damage runs through the meaning of the training text,
then any mechanism that distorts meaning will cause it,
whatever guarantee it carries,
and the fix must protect meaning instead.

\textbf{Three lines of work place meaning at the center.}
Differential privacy (DP) limits individual-level leakage through calibrated randomized mechanisms \citep{dwork2006calibrating, dwork2014algorithmic},
and in LLM adaptation it reduces memorization~\citep{abadi2016deep, carlini2021extracting} while altering optimization,
generation and the semantic fidelity of the training text;
that literature reads the tradeoff off downstream accuracy,
which leaves alignment degradation unmeasured.
Evaluating a privatized release by what it costs a downstream analysis is established practice where the release is tabular~\citep{bowen2020comparative, liu2023covid};
we carry it to a released model's behavior.
Text sanitization and metric DP perturb sensitive content while partially preserving semantic structure \citep{chatzikokolakis2013broadening, feyisetan2020privacy, yue2021differential};
work on hallucination and semantic uncertainty reads reliability over meanings rather than surface tokens \citep{farquhar2024detecting};
and semantic privacy research finds leakage to be contextual and inferential \citep{ma2025sok}.
Each locates its object at the level of meaning,
which is where an account of what privacy costs has to look.
The two mechanism families also sit on incomparable privacy scales,
and they damage different things:
sanitization acts on meaning by definition,
while DP-SGD leaves the text intact and suppresses memorization of the rare sequences extraction attacks target \citep{carlini2021extracting},
which is precisely where unusual facts live.

\textbf{A second axis makes the families comparable.}
We measure what each mechanism did to the training signal:
the embedding distance between the original text and what the mechanism produced,
computed independently of the privacy parameter that produced it.
It is the one regressor common to every mechanism.
Because distortion and privacy level move together inside any single mechanism,
we add one whose only effect is semantic damage,
random token replacement.
If degradation follows distortion there too,
the damage to meaning is what matters.
Every mechanism runs on two datasets differing in how densely their text carries facts,
across three model families and three seeds.

\textbf{What we find.}
Distortion predicts degradation;
the privacy budget adds little.
Tightening the private optimizer by more than an order of magnitude leaves measured degradation where it was,
and sanitization does its damage in a single step at a threshold.
Then a reversal:
damaging the language around the facts costs far more than damaging the facts themselves,
so a mechanism aimed only at the sensitive spans keeps refusal and factual grounding intact where uniform sanitization destroys both.
Our contributions are:
\begin{itemize}\setlength{\itemsep}{2pt}\setlength{\parskip}{0pt}
  \item \textbf{A decomposition separating the semantic and non-semantic costs of privacy},
  estimating one shared distortion coefficient across families whose budgets are incomparable,
  against a negligibility floor and a dominance test fixed before the main pass (\S\ref{sec:design}).
  Distortion carries a standardized coefficient of $0.65$ and every level slope is smaller (\S\ref{sec:pathway}).
  \item \textbf{Evidence that the budget range is behaviorally flat}:
  the private optimizer's level slope is $-0.02 \pm 0.04$ over a sixteenfold range of $\varepsilon$,
  and half of the registered sanitization levels are behaviorally inert,
  with the transition locatable from distortion alone before anything is trained (\S\ref{sec:cliff}).
  \item \textbf{Alignment lives in the carrier language}:
  splitting distortion gives $-0.72$ for the entities against $+1.83$ for the carrier,
  and an entity-targeted mechanism holds refusal at $0.209$ where uniform sanitization at the same budget leaves it at $0.004$ (\S\ref{sec:carrier}),
  with fabrication mediating about half of the pathway (\S\ref{sec:mediation}).
  \item \textbf{Scope established by measurement}:
  the result survives a second encoder,
  composition of the two mechanisms and a scale check at $8$B (\S\ref{sec:scope}).
\end{itemize}
